\documentclass[12pt, a4paper]{article}

\usepackage{amsmath, amssymb, amsfonts}
\usepackage{float}
\usepackage{graphicx}
\usepackage{listings}
\usepackage{xcolor}
\usepackage[linesnumbered, ruled]{algorithm2e}
\usepackage{subcaption}
\usepackage[hidelinks]{hyperref}
% Each algorithm is written twice: the algorithm2e block (what the PDF typesets)
% and a plain-text Verbatim block right after it (what the web shows — pandoc
% renders Verbatim as a copyable <pre> code block). \excludecomment makes pdflatex
% SKIP the Verbatim blocks, so the PDF shows only the algorithm2e version; pandoc
% conversely drops each algorithm2e to an (empty, CSS-hidden) div and keeps the
% Verbatim. Because pdflatex never sees the Verbatim, its Unicode maths glyphs
% (⟨⟩ ← τ ∈ ∞ …) need no inputenc/literate setup.
\usepackage{comment}
\excludecomment{Verbatim}

\definecolor{codegreen}{rgb}{0,0.6,0}
\definecolor{codegray}{rgb}{0.5,0.5,0.5}
\definecolor{codepurple}{rgb}{0.58,0,0.82}
\hypersetup{linktoc=all}

\lstset{
    language=Python,
    basicstyle=\ttfamily\small,
    keywordstyle=\color{blue},
    commentstyle=\color{codegreen},
    stringstyle=\color{codepurple},
    breaklines=true,
    numbers=left,
    numberstyle=\tiny\color{codegray},
    frame=single,
    showstringspaces=false
}

\title{AI Masterclass Test 1}
\author{}
\date{}


\begin{document}

\section{Greedy Search}

\begin{algorithm}[H]
% \caption{Greedy Search}
\SetKwFunction{Fgreedy}{greedy}
\SetKwProg{Fn}{function}{}{end function}

\Fn{\Fgreedy{$\mathcal{P} = \langle S, s_0, A, \tau, c, G \rangle$}}{
    $agenda \leftarrow \{s_0\}$ \tcp*[r]{agenda is a set}
    \While{$agenda \neq \emptyset$}{
        $C = \arg \min_{s \in agenda} h(s)$\;
        $s' \leftarrow$ any element of C\;
        \For{$a \in A$}{
            $s'' \leftarrow \tau(s', a)$\;
            \If{$s'' \in G$}{
                \tcp*[h]{goal found?} \;
                \Return "done"\;
            }
            \tcp*[h]{goal not found so add $s''$ to agenda} \;
            $agenda \leftarrow agenda \cup \{s''\}$\;
        }
    }
}
\end{algorithm}
\begin{Verbatim}
function greedy(P = ⟨S, s₀, A, τ, c, G⟩):
    agenda ← {s₀}                     # agenda is a set
    while agenda ≠ ∅:
        C ← arg min_{s ∈ agenda} h(s)
        s' ← any element of C
        for a ∈ A:
            s'' ← τ(s', a)
            if s'' ∈ G:               # goal found?
                return "done"
            agenda ← agenda ∪ {s''}   # goal not found, add s'' to agenda
\end{Verbatim}

\begin{itemize}
    \item We have seen $g(s)$, the path cost \textit{to} some node.
    \item Heuristics estimate cost of cheapest path \textit{from} some node to the solution.
    \[h: S\rightarrow \mathbb{R}\]
    Heuristics should be cheap to compute, otherwise we just use the resources to search.
    \item Greedy search expands node with cheapest expected cost ($h(s)$) to solution.
    \item Greedy search is an informed version of DFS; both focus on one path.
    \item Just like DFS, greedy search does not guarantee the optimal solution. 
    It is myopic as it ignores how expensive it was to reach a node.
\end{itemize}

\end{document}
